% ID: SEL-D0003-OA-E03
% ESTADO: cotejada. Fuente y fragmento cotejados de forma independiente; salida compilada revisada.
% Fuente: Selectividad/sel_2019_matematicas_aplicadas/Suplente Junio/Exámenes Matemáticas Aplicadas CCSS - A y B.pdf
% SHA-256: a59474e6895bd41034da53b7060969cbe8cbf7f8272f39fbf16821c95f06c84b
% Página: 1; opción A; ejercicio 3.
% Encabezado original: EJERCICIO 3. Solo enunciado; no contiene solución.

El 17\,\% de la población adulta de una ciudad sigue una dieta de
adelgazamiento y practica algún deporte regularmente. El 58\,\% ni sigue
una dieta de adelgazamiento ni hace deporte regularmente. Además, se sabe
que de los que hacen deporte regularmente, el 50\,\% hace dieta de
adelgazamiento. Se elige al azar un adulto de esa población.

\begin{enumerate}[label=\alph*)]
  \item \textbf{(1 punto)} \textquestiondown{}Cuál es la probabilidad de que
  siga una dieta de adelgazamiento o que practique deporte regularmente?
  \item \textbf{(1 punto)} Si el individuo elegido sigue una dieta de
  adelgazamiento, \textquestiondown{}cuál es la probabilidad de que practique
  deporte con regularidad?
  \item \textbf{(0.5 puntos)} \textquestiondown{}Son independientes los sucesos
  ``Seguir una dieta de adelgazamiento'' y ``Practicar algún deporte
  regularmente''?
\end{enumerate}
